% Fresh scope proof, 2026-10-22.  Retained proof bytes are unchanged.
\section{The actual hypotheses used by the retained weak-C2 leaves}
\label{sec:weak-c2-bare-source-audit}

This section proves the scope extension needed for the source join.
Several retained theorem statements were written on a complement of
earlier admissions.  Such a complement is not a mathematical
assumption in the coefficient calculations below.  The extension is
justified by checking the entire proofs, including the state normals
and the remote joins; it is not inferred from their historical
admission labels.

\paragraph{Common source hypotheses.}
Throughout this section use the original physical stationary C2
source, with \(d=1\), \(\lambda_{\rm event}=\beta=0\), its actual
contact, and all original strict gaps and nonnegative event
coefficients.  In particular \(n\geq0\), including equality.
Put
\[
\begin{gathered}
 z=aB,\quad y=tB,\quad H=hB,\quad L=lB,\quad
 l=1-h,\quad e=1-r,\\
 g=1-ae+haB,\quad j=1+aBr-l/t .
\end{gathered}
 \tag{BSR1}
\]
Here \(B,z,y\) are positive source infinitesimals, \(0<h,r<1\),
\(g,j>0\), and \(B-zyr>0\).
The full quarter-parent is the polynomial HJ2 in the following
source-assembly section.  All the retained identities cited here
are identities for that polynomial, not only for its value on a
contact graph.
A \emph{unit} below is positive, bounded, and has positive real
residue.  The source and testing fields are arbitrarily chosen,
separately labelled real closed extensions of \(\mathbb R\);
no algebraic-disjointness assumption is made.  They may be
two labelled copies of the same extension.  The source and
testing choices are independent in the universal quantifiers.
The positive-infinitesimal total-scale condition may be retained;
on every row with an unbounded original coefficient it follows
already from that coefficient.

\subsection{Regular A-opening exits}

On \(t,g,j\) unit and \(h\) infinitesimal, the four actual
source exits have exactly the following hypotheses:
\[
\begin{array}{c|l|l}
 &\text{additional source conditions}&\text{complete bodies}\\ \hline
 {\rm WA}&r,e,a\text{ unit}&
       \texttt{five\_root\_local},\ \texttt{five\_root\_remote}\\
 {\rm WE}&r,e\text{ unit},\ a\text{ infinitesimal}&
       \texttt{extension}\\
 {\rm SL}&r\text{ infinitesimal},\ a\text{ bounded}&
       \texttt{local\_small\_r},\ \texttt{remote\_small\_r}\\
 {\rm WC}&e\text{ infinitesimal},\ a\text{ unit}&
       \texttt{centered\_circuit\_closure},\
       \texttt{remote\_reduction}.
\end{array}
\tag{BSR2}
\]
The WA, WE and SL statements explicitly require no residual
membership premise.  Their source bounds use the original formulas,
and their complete five-root identities keep both state normals.
WE uses the exact monic \(a\)-division of
\((b_2-\alpha/B^2)/a\); its canceled denominator is a unit.
It requires no bound on \(h/a\).
SL uses \(nr\), \(Ah\) and the normalized roots
\(\sqrt n\,\xi,\sqrt A\,\eta,aR_2\); it requires no bound on
\(r/h\), \(a/r\), or a testing/source ratio.
Their printed homogeneous remote identities restore every
projective test.

For WC the precise extra scale bound is proved, rather than
assumed.  The exact original identity is
\[
 2r\Delta n=eQ-BK,\qquad
 K=(1+a)(h+t)(1+aB),
 \tag{BSR3}
\]
\[
 \begin{split}
 Q={}&1+B(1+t+ah+2ar+at)\\
    &+B^2\{a(h+t)+a^2[h(r+1)+t(1+r-r^2)]\}.
 \end{split}
\]
On the WC domain \(Q,K,\Delta\) are positive units.  Thus
\(B/e\) and \(n/e\) are bounded, including \(n=0\).
These are exactly BS1--BS3, not an unprinted assumption
about the relative rates of \(B,e,h\).
WC3 and WC13--WC18 prove the full scalar and normal coquotient
bounds; its reciprocal testing-ratio argument uses the same
physical \(n/m\) bound.  The finite three-circuit lift is used
only after these complete coefficient checks.  GR1--GR7 give
the remote parent, including the additional \(\sqrt m W\)
root.  None of these four exits calls a later weak-C2 row
or the present source join.

\subsection{Unit-\(e\) parts of R3 and R4}

The actual bodies \texttt{r3\_e\_unit\_all\_tests} and
\texttt{r4\_e\_unit\_all\_tests} apply on their explicit
row conditions together with \(e\) unit, without an earlier-exit
premise.  Indeed \(g>0\) first gives bounded \(a\).  In their
source calculations all denominators are units after replacing
the potentially divergent coefficients by
\[
 nr,\quad mg,\quad Cl\ \text{in R3},\qquad
 nr,\quad mg,\quad Dj\ \text{in R4}.
 \tag{BSR4}
\]
The full \(\xi\)-row is normalized by \(\sqrt{n+1/2}\).
R3 keeps \(\sqrt C\,\eta\); R4 uses the exact complete
centered-\(D\) identity RD3--RD6 and solves its \(l\)-unit
coefficient.  The only apparent \(a\)-pole is the same exact
WE monic quotient used in BSR2.  For \(g\) infinitesimal
the centered \(m\) excess is retained and the remainder
depends on bounded \(mg\), rather than on \(1/g\).
The strict two-by-two blocks R3E5--R3E6 and
R4E7--R4E8 include all these openings.
The remote proofs use only the printed positive lower bounds
and the complete selected-parent identities R4E14--R4E16.
Their source estimate \(nr\sim e/[2(1+a)]>0\)
proves that \(n=0\) is impossible on these particular domains.
It is not removed from a neighboring source branch.

\subsection{R5}

On bare R5, \(t,h,l,j\) are units and \(r\) is
infinitesimal.  Positivity of \(g\) gives bounded \(a\)
with real residue \(\rho\in[0,1]\).
Writing \(\chi=h_0\), \(\ell=1-\chi\), and \(\tau=t_0>\ell\),
the actual source calculation R5.1 gives
\[
 \begin{gathered}
 (nr)_0=\frac1{2(1+\rho)},\qquad
 A_0=\frac{\tau}{2\chi(1+\rho)},\\
 C_0=\frac{\chi+\tau-\rho\ell}{2\ell(1+\rho)},\qquad
 D_0=\frac1{2(1+\rho)(\tau-\ell)},\\
 (\alpha/B^2)_0
       =\frac{\chi+\tau-\rho\ell}{2(1+\rho)}.
 \end{gathered}
 \tag{BSR5}
\]
Every displayed pivot is positive.  For \(\rho<1\),
\(m/a^2\) is bounded; for \(\rho=1\), \(mg\sim1/4\)
and \(m\) is unbounded.  These alternatives are exhaustive,
and \(nr\) proves \(n=0\) impossible here.

The complete local proof uses \(\sqrt n\,\xi\), the
exact WE quotient when \(\rho=0\), and the centered
excess R5.2--R5.5 when \(\rho=1\).
For clarity, that last identity is
\[
 mM^2-(m-K)(M-gvW)^2
 =KM^2+2(m-K)g\,vWM-(m-K)g^2v^2W^2,
 \tag{BSR6}
\]
where \(M=vW-X+uvY-uZ\) and \(K>0\) is real.
It is an entire-event identity; it does not change the source
coefficients in any other packet.
The full \((\eta,R_1)\) determinant is
\(A_0(C_0+D_0+1/2)>0\).
The finite remote determinants R5.8 and R5.10, the
\(\rho=0\) extra root \(aW\), and R5.13--R5.17
cover all remaining states and projective tests.
Their entries use precisely BSR5, \(nr\), and \(mg\).
No earlier-admission condition or relative opening rate
occurs in this chain.  Therefore it proves bare R5.

\subsection{R6}

On bare R6, \(t,h,l,j,r\) are units and \(g\) is
infinitesimal.  The exact equality
\(a(e-H)=1-g\) excludes infinitesimal \(a\).
For bounded \(a\), \(e\) is a unit and \(a_0=1/e_0\).
R6.1 gives positive-unit \(n,A,C,D,\alpha/B^2,mg\).
The entire proof consists of BSR6 followed by the strict
local block R6.2, the full kernel determinant R6.4,
and the homogeneous matrices R6.5--R6.7.
Only these actual residues are used; no prior-complement
condition is needed.

For unbounded \(a\), put \(s=1/a\) and \(b=a^2B\).
The exact physical identity BSR3, in the original variables,
proves \(b\) bounded before any contact normalization.
The literal zero of \(n\) is retained in that inequality.
The positive bounded \(b\) is either a unit or infinitesimal:
\[
 \begin{array}{c|l}
 b\text{ unit}&\texttt{r6\_b\_unit\_all\_tests}
                       \quad({\rm RB1\!-\!RB24})\\
 b\text{ infinitesimal}&\texttt{r6\_b\_inf\_all\_tests}
                       \quad({\rm RI1\!-\!RI20}).
 \end{array}
 \tag{BSR7}
\]
Both proofs retain the actual centered \(m\) packet and,
where required, the whole actual \(n\) packet before forming
the bounded remainder.  RB5--RB9 give its full polynomial
inverse and actual contact.  RB13--RB18 and RI8--RI15
verify the raw scalar support and both complete first-normal
modules before applying the restricted CC1--CC6 conductor.
The infinitesimal-\(b\) proof verifies the additional
\(\varepsilon,\mathfrak j\) filtrations on the full coefficients;
it does not import an unrestricted contact-ideal inclusion.
Its strict leading block is RI19.
The adjacent unit-state charts retain the original \(n\)
packet when that is the needed pivot.

The remote inputs in the two dependency manifests are
selected full-parent identities: BE4, the BE infinity/zero
identity, R5.5, RT10, and R6.5--R6.7.
Their lower bounds are verified from the actual R6 source.
The C1 source theorem and a bounded-\(a\) R6 source
theorem are not transported to this domain.
Hence the three actual R6 bodies prove its entire bare
source.  \(n=0\) remains in the unit-\(b\) component;
where a positive residue excludes it, that exclusion is
proved by the original formula.

\subsection{R7 and its one genuine source exclusion}

Bare R7 has \(t\) unit and \(h,g\) infinitesimal, with
\(r,j\) not both infinitesimal.  This last condition
is used and must be kept.
For bounded \(a\), \(a_0=1/e_0\), and R7.1 gives
\[
 C_0=\frac{t_0-e_0}{2(1+e_0)},\qquad
 (\alpha/B^2)_0=\frac{t_0-e_0}{2e_0(1+e_0)},\qquad
 (Ah)_0=\frac{t_0}{2(1+e_0)}.
 \tag{BSR8}
\]
Here \(t_0\geq1\), \(0<e_0\leq1\).
The retained exclusion is exactly what prevents
\(t_0=e_0=1\); it makes the first two pivots strict.
The full proof \texttt{r7\_bounded\_a\_all\_tests}
uses the centered \(m\) identity and the exact
centered \(D\) event
\[
 L_D=uW+uvX-Y-vZ+jvZ .
 \tag{BSR9}
\]
Its five-root matrix identity R7.3 includes the zero
syzygy and both normals.  After extracting \(DL_D^2\),
the remaining coefficients involve bounded \(nr,Ah,mg,Dj\).
The appropriate hard rows are normalized jointly.
The remote cap lemma, full kernel matrices, and
R7.6--R7.12 use the same bounds, with no comparison
between the four divergent coefficients.

For unbounded \(a\), the actual source gives \(r\sim1\);
thus the retained \(r,j\) exclusion is automatic.
The body \texttt{r7\_unbounded\_a\_body} derives bounded
\(b=a^2B\), \(n=O(s^2)\), and bounded positive
\(mg,Ah,D(1-c)\) directly from the source formulas.
RU7 is the entire bounded remainder after all four
actual positive packets have been removed.
RU13--RU16 check its full raw support and normal modules;
RU20 verifies the coefficient filtrations before reusing
the restricted RI14--RI15 algorithm.
The complete strict block RU24 has determinant
\((t_0+2)/(4t_0)>0\), including \(t_0=1\).
RU26--RU30 and the selected BE identities supply the
remote and projective join.  No earlier-exit condition
is used.

The source join below proves exactly
\[
 r,g,h,j\text{ infinitesimal},\ t\text{ unit}
 \quad\Longrightarrow\quad C\text{ infinitesimal}.
 \tag{BSR10}
\]
Consequently the excluded R7 intersection cannot occur
on the directed strong-\(C\) input.  One must not extend
BSR8 across that intersection by deleting the condition.

\subsection{R8}

Bare R8 has \(t,g\) units and \(h,j\) infinitesimal;
the exact gap gives \(t_0=1\).
There are two exhaustive cases for \(r\).
For \(r\) a unit, \texttt{joint\_D\_R8\_closure}
gives the complete local proof JD1--JD18.
The centered-\(D\) identity JD4 removes the whole
divergent rank-one block, including both normal terms.
It then splits the actual source into bounded \(a\)
with \(e\) unit, bounded \(a\) with \(e\) infinitesimal,
and unbounded \(a\).
The latter uses the positive source quantity
\(\rho=n/e^2\): its noninfinitesimal part has the
strict weighted pivot, its positive infinitesimal
part has the exact paired-\(D\) thin lift, and
\(n=0\) has the physical map described below.
JD14--JD17 verify the new \(D\)-dependent coquotients;
no inverse of \(j\) remains in their bounded arrays.
This is an algebraic comparison with a formal source,
not application of a physical C1 theorem at negative beta.

For \(r\) infinitesimal, the independent body
\texttt{r8\_small\_r\_local} derives
\[
 (nr)_0=(Ah)_0=(Dj)_0=\frac1{2(1+a_0)},\qquad
 0\leq a_0<1 .
 \tag{BSR11}
\]
It jointly normalizes \(\sqrt n\,\xi\), \(\sqrt A\,\eta\),
and \(\sqrt D\,L_D\), using the exact matrix syzygy
RD3--RD6.  RD7 checks every full remaining coefficient,
including the monic \(a\)-quotient at \(a_0=0\).
Its five strict residue pivots are printed in RD9--RD11.
No relative rate among \(r,h,j\) is required.
JD19--JD24 and the exact homogeneous selected-parent
identities cover the remote complement for both cases
of \(r\).  Thus their union proves bare R8; it uses no
historical residual-exclusion premise.

\subsection{R9 and R10}

Bare R9 has \(t,g,j\) units and \(h,e,a\)
infinitesimal.  Its exact source denominator has unit
residue \(2(t_0-1)>0\).
The R9 source appendix proves the cancellations
\[
 \lambda=\frac{m-a^2q}{ta^2B}\sim\frac12,\qquad
 c_F\sim\frac1{2t_0},\qquad
 C\sim Ah\sim\frac{t_0}2,\qquad
 D\sim\frac1{2(t_0-1)}.
 \tag{BSR12}
\]
The body \texttt{reciprocal\_contact\_closure},
RC1--RC13, contains the entire scalar, both complete
first normals, and the hard block.  Its coupled circuit
removes the apparent unbounded testing-ratio coefficient;
the final homogeneous roots have no such inverse.
The R13 frontier's R9 extension of GR1--GR7 uses
the actual bounds \(B/e,n/e\) and \(m/(a^2e)\),
the full additional \(\sqrt m W\) row, and the
printed homogeneous endpoints.
All these estimates are on the explicit bare source,
including \(n=0\).

Bare R10 has \(t,g,j\) units, \(h,e\) infinitesimal,
and \(a\) unbounded.  The R10 source-and-weighted proof puts
\(s=1/a,w=ae,b=a^2B\) and proves
\[
 b/w=O(1),\qquad n=O(s^2w),\qquad
 m/w\sim\frac1{2(1-w_0)},\qquad w_0<1 .
 \tag{BSR13}
\]
The bound on \(b/w\) comes from the original inequality
BSR3 before any normalized circuit is introduced.
The same source proof gives the full remote parent,
retaining bounded \(m\) as a whole packet where needed.
The actual all-tests join has three local source cases:
\(\rho=n/e^2\) noninfinitesimal, \(n=0\), and
\(0<\rho\) infinitesimal.
The first is the complete weighted identity R10.9
and strict source Schur cases R10.11--R10.12.
The last is the whole thin-parent identity TL4--TL10
together with \texttt{r10\_unit\_W\_join}.
The latter proves the exact critical deficit divided
by \(m\) equals \(s\) times a bounded source coefficient;
it includes every bounded normalized \(Z\), including zero.
The maximal-coordinate state partition supplies exactly
that chart and its complementary \(Z\)-unit chart.
The restricted CA/CS/SE conductors are used only after
the full raw supports, both normal arrays, and the
required scaled divisibilities have been checked.
No bound on \(e/B\), \(h/e\), or an independent
testing/source ratio is inserted.

\paragraph{The literal zero boundary in R8 and R10.}
The only physical C1 source import just used is
at literal \(n=0\).
The exact OP4 map is
\[
 x_1=z,\quad \eta_1=t,\quad \zeta_1=1/a,\quad
 r_1=h,\quad \theta_1=1-r .
 \tag{BSR14}
\]
It gives the identical four contact state coordinates.
OP2 and OP5 prove that the formal beta is exactly zero
on this boundary.  The original C2 parent is therefore
the actual C1 parent, up to its allowed positive
coefficient normalization.
In the unbounded-\(a\) cases where this import is used,
\[
 k_1=\sqrt z,\ \zeta_1,\ r_1\text{ are infinitesimal},
 \qquad \eta_1=t\text{ is a unit},\qquad
 e_{1,\mathrm{scale}}^2=t\,a^2B
 \text{ is positive bounded}.
 \tag{BSR15}
\]
If the last scale is infinitesimal, apply the actual
theorem \texttt{small\_e\_all\_tests}.
Otherwise apply \texttt{centered\_a\_bounded\_delta}.
The former covers every positive native \(r_1\);
the latter covers the present infinitesimal native \(r_1\).
Both statements retain every original positive opening
and literal \(\beta_1=0\), and neither has an
earlier-complement premise.
Their full centered-parent and both-normal proofs,
outer/state lemmas, and homogeneous endpoints are
the stated imports.  No whole-C1 routing theorem
or later weak-C2 result is a dependency.
For \(n>0\), OP5 gives negative formal beta;
BSR14 is then never used as a physical-source theorem.

\paragraph{Conclusion of the scope proof.}
Every use of a retained source theorem has now been
matched to its explicit algebraic hypotheses.
Every use outside its old source label is an exact
selected-parent identity or a restricted conductor
whose complete coefficient hypotheses have been proved
in the importing body.
In particular all full hard blocks and both state-normal
rows remain in the respective identities.
The only retained source exclusion needed by the
ten-row join is the simultaneous \(r,j\) condition
in R7, and BSR10 dispatches it explicitly.
Thus the listed bodies prove the bare-source versions
needed for HJ, without a circular appeal to the
union being assembled.  Every file and the exact
imported locations are pinned in the accompanying
used-input manifest.
